\section{总结与延伸}

这堂课用四个主题回答“Beyond LLMs”意味着什么。Emergent abilities 提醒我们区分真实能力变化与 metric artifact；Intermediate-Guided Reasoning 把一次生成扩展为 chain、tree、decomposition 或 program execution；BabyLM 用受限数据预算追问 sample efficiency；Agents 则把模型放进能行动、记忆、协作和纠错的系统。

\subsection{统一框架：能力来自模型、过程与闭环}

前后半场可以统一为三层：model 提供概率化表示与生成，inference procedure 组织 intermediate states，system loop 把结果接入 environment。规模更大不自动带来忠实推理，推理轨迹更长不自动带来正确行动，工具更多也不自动带来可靠系统。每一层都需要独立指标和失败分析。

\begin{importantbox}{本讲的八条可执行结论}
\begin{enumerate}
  \item 报告 emergence 时同时检查 metric、prompt、task distribution 与连续替代指标；
  \item 为 reasoning method 区分 final-answer accuracy、step validity 与 compute cost；
  \item 精确计算优先交给 interpreter，并验证程序与问题语义的一致性；
  \item 比较小模型时同时报告 data budget、architecture、objective 与 curriculum；
  \item 把 foundation model 与 agent system 分开评测；
  \item 给每个真实动作定义 permissions、success predicate 与 confirmation gate；
  \item Multi-agent 使用结构化 task state、有限重试和 evidence aggregation；
  \item 长时运行必须依赖 external feedback、audit、sandbox 与 recovery。
\end{enumerate}
\end{importantbox}

\subsection{如何验证一个 Agent Claim}

面对“agent 能订票、通过考试、操作任意软件”之类陈述，可按五步检查：先固定任务与版本，再公开初始状态和权限；随后运行多次并保存完整 trajectories；用环境 postcondition 而非模型自评判定成功；最后报告失败类别、人工介入、成本和不可逆事件。这样才能把 demo 转成可比较的 evidence。

\begin{warningbox}{不要把 2023 年课堂愿景写成 2026 年产品事实}
本讲记录的是 2023-11-28 的课程内容。MultiOn、AutoGPT、Action API 与 LLM OS 均在特定历史语境中出现；讲义保留其教学价值，但不据此断言今天的产品能力、公司状态或行业标准。
\end{warningbox}

\subsection{拓展阅读与最终小结}

\enlargethispage{5\baselineskip}

建议沿三条路线继续：第一，阅读 emergent abilities 与 metric-choice 争论，比较离散 exact-match 和连续损失；第二，阅读 Chain-of-Thought、Tree of Thoughts、PAL、Program of Thoughts 与 BabyLM Challenge，理解推理过程和样本效率；第三，阅读 Lilian Weng 的 LLM-powered autonomous agents 综述，并结合现代 agent evaluation、computer-use safety 与 sandbox 设计重新审视课堂架构。阅读时应优先论文、官方项目页和可复现实验，不以产品 demo 代替证据。

真正的“Beyond LLMs”不是抛弃 LLM，而是不再把一次文本生成当作能力终点。更可信的系统会选择合适的 intermediate representation，用工具获得可验证结果，用 memory 保存有 provenance 的状态，用 multi-agent 并行 bounded tasks，并在 permissions、feedback、reflection、sandboxing 与 human oversight 下执行现实动作。

\end{document}
